% Appendix D --- orientation, and the orders generated by invariant cones. Moved out of section 5
% (Q-0290, the author's decision of 2026-10-04 to shorten section 5 by moving material to appendices,
% following the blind review records/affine/REVIEW-paper-section-5.md, flags S21 and the length
% judgement on 5.3). Statements moved whole with their `% shared with blueprint` markers, verbatim;
% the main text keeps a sentence naming each and pointing here.
%   - prop:glplus-families (was in 5.2): the paragraph "Orientation" of 5.2 states its use.
%   - lem:spectral-cones (was in 5.3): Ky Fan's minimum principle and the log-majorization of
%     products are proved here, not cited: no matrix-analysis source is held by the librarian
%     (`library search --corpus` for Bhatia, Horn--Johnson, Marshall--Olkin, Ky Fan: none, 2026-10-04).
%   - the proof of prop:invariant-orders, whose statement stays in 5.4.
%   - thm:cone-orders with its proof (second pass). The review moves its clauses (2)-(3) out; the node
%     is not split (a blueprint change, owed to the session after the Lean skeleton), so it moves
%     whole and 5.4 keeps a paragraph stating all four clauses in words.
% Q-0309 (fix group 2; records/affine/move-q0309-flow.py): this file was appendix-orders.tex. The
% statement of thm:cone-orders moved to 5.4 (review flag 35); its proof, lem:spectral-cones and the
% proof of prop:invariant-orders moved to the proof appendix (appendix-proofs.tex), so that the
% appendices are the proofs by section first, then the supplementary results (flag 39). What is left
% is the orientation-preserving group, Appendix C.

\section{The orientation-preserving group}\label{app:orientation}

Covariance under the orientation-preserving matrices $\GL^+(d)$ alone is a weaker requirement than
covariance under $\GL(d)$: it leaves out the reflections. For $d \ge 2$ it gives the same families. An
order invariant under $\GL^+(d)$ is invariant under $\GL(d)$, and the $(\GL^+(d), W)$-cascade families of
kernels and of operators are the $(\GL(d), W)$-cascade families (Proposition~\ref{prop:glplus-families}).
So the orientation-preserving group adds no row to Table~\ref{tab:landscape}, as \S\ref{sec:scale-shape}
states. On the line the reflection matters, and clause (3) gives a family of $\GL^+(1)$ that is not one
of $\GL(1)$.

% shared with blueprint prop:glplus-families+a4ce8d6baf6e
\begin{proposition}[covariance under $\GL^+(d)$ is covariance under $\GL(d)$]\label{prop:glplus-families}
Let $d \ge 1$ and $\GL^+(d) := \{A \in \GL(d) : \det A > 0\}$.
\begin{enumerate}
\item \emph{(The orders.)} $\GL^+(d)$ is a closed subgroup of $\GL(d)$ containing the dilations. For
all $\Sigma, \Sigma' \in \PSD(d)$ there is a $B \in \GL(d)$ with $\det B = -1$, $B\Sigma B^{\tp} = \Sigma$
and $B\Sigma'B^{\tp} = \Sigma'$. Consequently every relation on $\PSD(d)$ that is invariant under
$\Sigma \mapsto A\Sigma A^{\tp}$ for every $A \in \GL^+(d)$ is invariant under it for every
$A \in \GL(d)$, and $(\GL^+(d), W)$ is a wedge datum (Definition~\ref{def:gw-cascade-family}) if and only if
$(\GL(d), W)$ is one.
\item \emph{(The families, $d \ge 2$.)} Let $d \ge 2$ and let $(\GL^+(d), W)$ be a wedge datum. A
family of probability measures indexed by $P_W$ is a $(\GL^+(d), W)$-cascade family of kernels if and
only if it is a $(\GL(d), W)$-cascade family of kernels; and a family of operators indexed by $P_W$ is
a $(\GL^+(d), W)$-cascade family of operators if and only if it is a $(\GL(d), W)$-cascade family of
operators. In particular, for $W$ the Loewner order, the $(\GL^+(d), W)$-cascade families of kernels
are the affine Gaussian scale spaces \ref{eq:affine-gaussian}, $b > 0$.
\item \emph{(The line.)} Let $d = 1$ and let $W$ be the Loewner order, which on $\PSD(1) = [0,\infty)$
is the usual order. The point masses $\mu_{s,t} := \delta_{\sqrt t - \sqrt s}$, $0 \le s \le t$, form a
$(\GL^+(1), W)$-cascade family of kernels that is not a $(\GL(1), W)$-cascade family.
\end{enumerate}
\end{proposition}

So for $d \ge 2$ every statement made in \S\ref{sec:landscape} for $(\GL(d), W)$ holds verbatim for
$(\GL^+(d), W)$. Two facts carry the proposition. An orientation-reversing matrix acts on a given pair
as an orientation-preserving one does, through the joint stabilizer of clause (1). And a law
invariant under $\mathrm{SO}(d)$ is invariant under $\Orth(d)$ when $d \ge 2$, because $\mathrm{SO}(d)$
is transitive on spheres. The second fact fails on the line, and clause (3) shows that the conclusion
fails with it (Remark~\ref{rem:landscape-line}). A group that
contains $\mathrm{SO}(d)$ but not $\Orth(d)$ and is not transitive on $\PD(d)$, such as the
orientation-preserving similarities, is part of problem~\ref{prob:continuous-fields} of
\S\ref{sec:open-problems}.

% proof of prop:glplus-families
\begin{proof}[Sketch]
\emph{(1).} Put $T := \Sigma + \Sigma'$, and for $u \ne 0$ let
$H_u := I - 2\langle u, u\rangle^{-1}uu^{\tp}$, the reflection in $u^\perp$, of determinant $-1$; it
commutes with every symmetric matrix of which $u$ is an eigenvector. If $T$ is singular, a $u \ne 0$
with $Tu = 0$ has $\Sigma u = \Sigma'u = 0$, and $B := H_u$ fixes both. Otherwise let $R := T^{1/2}$
and $u$ an eigenvector of $M := R^{-1}\Sigma R^{-1}$, which is also one of $R^{-1}\Sigma'R^{-1} = I - M$;
then $B := RH_uR^{-1}$ fixes both and has $\det B = -1$. For $\det A < 0$, the matrix
$AB \in \GL^+(d)$ acts on the pair as $A$ does, which
carries invariance from $\GL^+(d)$ to $\GL(d)$. Both groups are closed and contain the dilations, and
the other conditions of a wedge datum do not involve the group, so $(\GL^+(d), W)$ and $(\GL(d), W)$
are wedge data together.

\emph{(2).} The conditions (K1), (K2) and (K4) of Definition~\ref{def:gw-cascade-family} do not
involve the group, so only (K3) for $\det A < 0$ is to be shown. For $d \ge 2$ the orthogonal
matrices of determinant $1$ act transitively on each sphere, so (K3) at $(0, I)$ makes $\mu_I$
rotation invariant, $\widehat\mu_I(\omega) = F(|\omega|^2)$. Then (K3) with $A = \Sigma^{1/2}$ and (K2) give
$\widehat\mu_\Sigma(\omega) = F(\langle\omega,\Sigma\omega\rangle)$ on all of $\PSD(d)$. Along the ray of
$I$ put $f(s,t) := \widehat\mu_{sI,tI}(\omega)$. By (K2) it is continuous, and by (K1) $f(t,t) = 1$ and
$f(r,t) = f(s,t)\,f(r,s)$. At a first zero $t^* > 0$ of $f(0,\cdot)$, continuity at $(t^*,t^*)$ gives an
$s < t^*$ with $f(s,t^*) \ne 0$, and then $f(0,t^*) = f(s,t^*)\,f(0,s) \ne 0$, a contradiction. So no
transform vanishes. The increments are then the quotients $\widehat\mu_{\Sigma'}/\widehat\mu_\Sigma$,
and (K3) follows for every $A \in \GL(d)$ by the uniqueness of the transform. For operators,
Lemma~\ref{lem:operators-to-kernels}(1), (2) and (4) carry the statement over. For the Loewner order
the families of kernels are the affine cascade families of
Definition~\ref{def:affine-cascade-family}, which Theorem~\ref{thm:affine-rigidity} identifies with
the affine Gaussian scale spaces.

\emph{(3).} $\GL^+(1) = (0, \infty)$ acts by $a \cdot s = a^2s$, and the point masses satisfy (K1)
to (K4) by inspection. The matrix $-1 \in \GL(1)$ fixes every $s$, while
$(-1)_*\delta_1 = \delta_{-1} \ne \mu_{0,1}$.
\end{proof}
